\section{Lower-dimensional estimates}\label{sec:low}

This section supplies lower-dimensional estimates and excludes extremal
configurations in which both shear costs vanish. Its two external inputs
are the wave-envelope estimate of Guth, Wang, and Zhang and the planar
Furstenberg theorem of Ren and Wang. We prove the Hilbert-valued,
weighted, and probabilistic consequences used here.

The wave-envelope estimate gives a kinetic bound for the base rays.
A separate scalar-channel estimate controls angular memory, which will
also be used with the canonical names in Section~\ref{names:section}.
The kinetic bound excludes hybrid zero costs when the positive hybrid
exponent is strictly larger than the classical exponent. The planar incidence input,
also used in the later projection arguments, excludes classical zero
costs when $\Gamma>\eta/2$. These exclusions leave the positive-cost
configurations for Sections~\ref{names:section}--\ref{act:section}.
The polynomial cutoff convention, the order of
limits, and the meaning of negligible errors are those of the cylinder
framework. In particular, the number of masks is fixed during each inner
scale limit.
Constants depending on that number are permitted; an ordinary power of
the scale depending on that number is not.

\subsection{The wave-envelope estimate and a kinetic bound}

The scalar wave-envelope estimate below is
\cite[Theorem~1.3]{GuthWangZhang2020}, in the smooth-cutoff form used here.
The following proof derives its Hilbert-valued and weighted extensions.

\begin{theorem}[Wave-envelope estimate and extensions]
\label{low:gwz}
Let $L\geq 2$, and let the Fourier transform of $F$ be supported in an
$O(L^{-1})$-neighborhood of a fixed compact annular patch of the cone in
$\R^3$. Use a smooth, boundedly overlapping decomposition into angular
sectors $\vartheta$ of width $L^{-1/2}$. For each dyadic
$L^{-1/2}\lesssim R\lesssim 1$, let $\tau$ range over a boundedly
overlapping cover by angular caps of width $R$. Let $U_{\tau,L}$ be the
associated envelope box, with dimensions $L$, $LR$, and $LR^2$ in the
adapted ray frame, and tile space by its translates. Then, for every
$\eps>0$,
\begin{equation}
 \norm{F}_{L^4}^4
 \lesssim_{\eps} L^{\eps}
 \sum_{L^{-1/2}\lesssim R\lesssim 1}
 \sum_{\tau:\,|\tau|\sim R}
 \sum_{U\parallel U_{\tau,L}} |U|^{-1}
 \left(\int_U\sum_{\vartheta\subset O(\tau)}
                    \norm{F_\vartheta}^2\right)^2.
 \tag{GWZ}\label{eq:GWZ}
\end{equation}
The assertion holds for functions taking values in a separable Hilbert
space, with a constant independent of that space. Fixed thickness
constants, bounded enlargements of tiles, and adapted rapidly decreasing
tile weights are allowed.
\end{theorem}

\begin{proof}
For the scalar case, apply the cited Theorem 1.3.
Here the enlarged cap in the displayed formula accommodates the bounded
overlap of the smooth partition. For completeness, the Hilbert-valued
extension introduces no dimension-dependent loss. In a finite-dimensional
subspace apply the scalar theorem to $\ip{F}{G}$, where $G$ is a standard
complex Gaussian vector. Its fourth moment is a fixed constant times
$\norm{F}^4$. Moreover,
\[
 \E\bigl(|\ip{v}{G}|^2|\ip{w}{G}|^2\bigr)
 =\norm{v}^2\norm{w}^2+|\ip{v}{w}|^2
 \leq 2\norm{v}^2\norm{w}^2.
\]
Average the scalar inequality and use this bound inside each double
integral on the right. Orthogonal finite-dimensional projections and
monotone convergence give the separable case. The weighted versions
follow by summing translates with rapidly decreasing coefficients;
changing an adapted smooth cutoff amounts to convolution on its dual
scale and gives the same summable enlargement.
For a scalar formulation localized by a decaying spatial weight, see also
\cite[Theorem~1.5]{GuthWangZhang2020}.
\end{proof}

We will use a fixed linear image and dilation of this theorem. At
frequencies of size $L$, consider
\[
 \rho=\frac{\xi^2}{4\zeta},\qquad \zeta\asymp L,
 \qquad \left|\frac{\xi}{2\zeta}\right|\lesssim 1,
\]
with thickness $O(1)$. The change
$(\rho,\xi,\zeta)\mapsto(\rho+\zeta,\rho-\zeta,\xi)$ identifies this
quadratic cone with a circular cone. A cap centered at the parameter
$x_0=-\xi/(2\zeta)$ has physical envelope of bounded length in direction
$(1,x_0,x_0^2)$ and transverse widths $R,R^2$ in the coordinates
\[
 b-x_0t,\qquad y-2x_0b+x_0^2t.
\]
Its volume is comparable to $R^3$. These coordinates also specify the
orientations in the following positive-measure consequence.

\begin{lemma}[Kinetic estimate for parabolic rays]\label{low:kinetic}
Let $0<\delta<1/2$. Let $\mu$ be a finite positive measure on rays
$z=(X,B,Y)$, of total mass $E$, supported where $|X|\lesssim1$. Assume,
for every center and $\delta^2\lesssim R\lesssim1$,
\[
 \mu\bigl\{|X-x|\leq R,\ |B-b|\leq R,
       \ |Y-y-2x(B-b)|\leq R^2\bigr\}\leq\kappa R^3.
\]
For each of a $\delta$-separated set of times in a fixed bounded interval,
let $D$ range over squares of side comparable to $\delta$. Assume bounded
overlap also after a fixed enlargement. With
$Q_z(t)=(B+tX,Y+tX^2)$, one has
\begin{equation}
 \delta^{-1}\sum_{t,D}\mu\{Q_z(t)\in D\}^2
 \lesssim_{\eps}\delta^{-\eps}\kappa E.
 \tag{kin}\label{eq:kin}
\end{equation}
\end{lemma}

\begin{proof}
Set $L=\delta^{-4}$ and form the $L^2(\mu)$-valued function
\[
\begin{split}
 F(t,b',y';z)={}&\delta^{-3}\chi(t)
 \phi\left(\frac{b'-B-tX}{\delta^2}\right)\\
 &\quad\cdot
 \phi\left(\frac{y'-Y-tX^2-2X(b'-B-tX)}{\delta^4}\right)
 e^{iL(X^2t-2Xb'+y')}.
\end{split}
\]
Choose Schwartz functions $\chi,\phi$ with sufficiently small compact
Fourier support. Require $\chi$ to be bounded away from zero on the time
interval under consideration and $\phi$ to be bounded away from zero
near zero. The three packet coordinates show that
\[
 \zeta\asymp L,\qquad -\frac{\xi}{2\zeta}=X+O(\delta^2),\qquad
 \rho+X\xi+X^2\zeta=O(1).
\]
Consequently
$\rho-\xi^2/(4\zeta)=O(1)$, so Theorem~\ref{low:gwz} applies in the
scaled coordinates just described. Spatial integration in the two packet
coordinates gives $\norm{F}_{L^2}^2\lesssim E$.

A fine angular projection has a uniformly smooth multiplier in these
rescaled frequency coordinates. Thus its packet components vanish unless
$X$ belongs to a bounded enlargement of that fine angular sector, and
they retain arbitrary rapid decrease in the same spatial packet
coordinates. We claim that for an $R$-cap and any of its tiles,
\[
 \int_U\sum_{\vartheta\subset O(\tau)}
             \norm{F_\vartheta}^2\lesssim\kappa R^3.
\]
Only $X$ within $O(R)$ of the cap center $x_0$ occurs. A ray meeting a tile
near time $t_0$ backflows into horizontal width $O((1+|t_0|)R)$ and
tilted vertical width $O((1+|t_0|)R^2)$. Indeed its tilted vertical velocity
relative to the central ray is $(X-x_0)^2=O(R^2)$. These enlarged boxes
can be covered by a polynomial number in $1+|t_0|$ of the boxes in the
hypothesis. The time decay of $\chi$ absorbs that polynomial. For packets
not meeting the tile, sum spatial shells in packet units: their energy
decreases faster than any prescribed power, whereas the number of
regularity boxes needed in the shell grows only polynomially. Integrate
the normalized spatial density of each packet first, and then integrate
against $\mu$. This proves the claim, including weighted tiles.

At a fixed cap scale the sum of these tile integrals is $O(E)$ by overlap.
Since $|U|\asymp R^3$, Equation~\eqref{eq:GWZ} gives
\[
 \int\norm{F(t,b',y')}^4\,dt\,db'\,dy'
 \lesssim_{\eps}\delta^{-\eps}\kappa E.
\]
For a square counted at time $t$, throughout a sufficiently short interval
of length comparable to $\delta$ the packet cores of all counted rays
lie in a fixed enlargement of the square. The integral of $\norm{F}^2$
over that enlargement is therefore at least a constant times the counted
mass. Cauchy--Schwarz in a spatial region of area $O(\delta^2)$, followed
by integration over the short time interval, contributes at least
$c\delta^{-1}\mu\{Q_z(t)\in D\}^2$. Bounded overlap in space and time
proves Equation~\eqref{eq:kin}.
\end{proof}

\subsection{The scalar channel}

In this subsection the base variables, if present, are retained only as
exact indices. The observable variables are the normal velocity and
position. In the hybrid model the fiber indexed by $\sigma\in[1,2]$
evolves by the Fourier multiplier $e^{-itHp^2/(2\sigma)}$, so that its
velocity is $N=Hp/\sigma$. Write $d_\ell=(H/\ell)^{1/2}$ and
$q_\ell=(H\ell)^{1/2}$. In the classical model $d_\ell=0$.

\begin{proposition}[Scalar-channel estimate]\label{low:scalar}
Suppose the root analyzed measure has mass $e$ and satisfies
\[
 \mu_0\{|N-n|\leq f,\ |D-b|\leq f\}\leq b_0 f^{1-\eta}
 \qquad(f>0,\ f\geq d_1).
\]
Here the subscript $1$ means time length one. At all terminal intervals
of length $\tau$, make disjoint fractional assignments to tests of normal
velocity width $f_\lambda$ and current-position width $\tau f_\lambda$,
where $f_\lambda\geq d_\tau$. The centers are constant on each test;
no shear depending on an exact index is used in these tests. With any
fixed finite number of admissible masks,
\begin{equation}
 \sum_\lambda \frac{m_\lambda^2}{f_\lambda^{1-\eta}}
 \leq \tau^{-\eta}s^{-o(1)}b_0e.
 \tag{SC}\label{eq:SC}
\end{equation}
This holds along polynomial sequences for which $\tau^{-1}$,
$b_0/e,e/b_0$, the phase density and support bounds, and $H,H^{-1}$
when applicable are bounded by fixed powers of $s^{-1}$.
\end{proposition}

\begin{proof}
We first prove the free estimate. The density bound controls the mass of
a small test by $s^{-C}ef^2$; for very large $f$ use total mass.
Consequently widths and reciprocal widths outside a sufficiently large
polynomial range contribute an arbitrarily small power, by summing
assigned mass. Split the remaining widths into $O(1+\log(1/s))$ dyadic
classes. Normal dilation reduces each class to $f=1$; the change in the
weight is canceled by the inverse change in the regularity constant.
In the hybrid case admissibility now gives $H\leq\tau$.
Each test is contained in boundedly many lattice cells of dimensions
$1\times\tau$. Splitting assignments between these cells and combining
assignments in the same cell reduces their square sum to the square sum
of the cell masses, up to an absolute constant.

For classical free evolution, express the latter square sum as a
collision count against $\mu_0\otimes\mu_0$. Pairs with velocity
separation comparable to $r$, $\tau\leq r\lesssim1$, can occupy a common
cell at only $O(r^{-1})$ sample times. Such a collision also requires
initial position separation $O(r)$. For each first ray the regularity
hypothesis bounds the mass of the second rays in question by
$Cb_0r^{1-\eta}$. The contribution of this separation scale is at most
$Cb_0er^{-\eta}$. Pairs separated by less than $\tau$ use at most
$\tau^{-1}$ times and lie initially in a box of width $O(\tau)$; their
contribution is at most $Cb_0e\tau^{-\eta}$. Sum the logarithmically many
scales. Velocities in a common unit cell differ by only a fixed constant,
so there are no additional separation scales.

We give the hybrid proof in detail. If $\tau=s^{o(1)}$, packet locality
and contraction bound the mass of each output cell at one time by
$s^{-o(1)}b_0$. Its input predecessors lie in a root regularity box of
width $s^{-o(1)}$, since $d_\tau\leq1$. Sum cell masses and the
$\tau^{-1}=s^{-o(1)}$ times. Hence suppose $\tau$ has positive limiting
depth, and use $\tau$ as the inner scale. All cutoff powers remain finite.
Fix an arbitrarily small $\eps>0$. Every loss below is
$\tau^{-C\eps}$ with a fixed $C$; tails outside $\tau^{-\eps}$-enlarged
adapted neighborhoods are made negligible before $\eps$ is sent to zero.

Partition the input by smooth compactly supported unit Fourier-velocity
caps with bounded overlap. Their energies $e_{\rm cap}$ sum to $O(e)$.
The Gaussian final analyzer sees, in any unit output velocity window,
only caps within distance $O(\tau^{-\eps})$, modulo negligible operator
norm. Thus Cauchy--Schwarz loses only $\tau^{-O(\eps)}$ on reducing the
square sum to the contributions of individual caps. Gaussian decay
justifies summing this error over all relevant caps and cells; the
polynomial support convention also permits restricting to polynomially
many significant caps. Boost each cap to $|N|\lesssim1$, and call its
input $g$ and its free evolution $u(t,x)$.

Plancherel in the analyzer's velocity variable identifies its spatial
marginal with a Gaussian average of $\norm{u(T,\cdot)}^2$ of width
$q_\tau\leq\tau$. For any $t\in[T,T+\tau]$, the band's spatial mass
on an interval transfers, in either time direction, into its
$O(\tau^{1-\eps})$-neighborhood with negligible $L^2$ error. To verify
this, insert a smooth Fourier plateau on the band. Away from velocities
in its support, integration by parts in its propagation kernel gives a
rapidly decreasing Schur tail. Its wavelength is $O(H)\leq O(\tau)$,
which is smaller than the allowed neighborhood. The retained operator is
uniformly bounded on $L^2$.

At each $t$, enlarged spatial cells overlap $O(\tau^{-\eps})$ times.
Average the preceding comparison in each terminal interval and expand
the square of a cell integral. Up to negligible errors and
$\tau^{-O(\eps)}$, the cell square sum is bounded by
\begin{equation}
 \tau^{-1}\int_{|\Delta|\lesssim\tau^{1-\eps}}
 \int_0^1\int_\R
       \norm{u(t,x)}^2\norm{u(t,x+\Delta)}^2\,dx\,dt\,d\Delta.
 \tag{pair}\label{eq:pair}
\end{equation}
Dropping the velocity cell after fixing a cap has only the already
allowed multiplicity. A boost merely translates the spatial grids at
each time. Errors can first be summed over polynomially many cells
covering significant support; the remaining mass is negligible.

For the Fourier calculations, first fix two exact indices, with parameters
$\sigma,\sigma'$, and view the integrand as the squared norm of the
tensor product of their fields. We integrate the resulting inequalities
over this pair of indices before applying input regularity. In particular,
all energies in the regularity bounds below include integration over their
exact indices. Decompose
Fourier-velocity pairs into Whitney scales
\[
 r_0=\max(\tau,\sqrt H)\lesssim r\lesssim1.
\]
At scale $r$, use products of smooth multipliers $a_I,b_J$ on intervals
of width $O(r)$, with center separation $O(r)$ and boundedly many
partners per interval. Require separation $\gtrsim r$, except for the
nearest pairs at $r_0$. One explicit construction starts with smooth
lattice partitions $\psi_i^r$, supported on $|N-ir|\leq r$, and the
near-pair symbol $\sum_{|i-j|\leq C}\psi_i^r(N_1)\psi_j^r(N_2)$.
Telescope these symbols through dyadic scales. Expand a difference using
the partitions at both adjacent scales in both variables. If $C$ is a
sufficiently large fixed constant, every surviving tensor term is
separated by a constant times $r$ and only boundedly many terms occur
locally. The last near-pair symbol is the remainder at $r_0$.

Insert a Schwartz time cutoff with compact Fourier support and modulus
bounded below on $[0,1]$. At fixed $\Delta$ and fixed indices, the
spacetime Fourier supports of products at any one Whitney scale have
bounded overlap. After multiplying frequency coordinates by $H$, these
supports satisfy
\[
 \sigma N_1+\sigma'N_2,\qquad
 \sigma N_1^2+\sigma'N_2^2+O(H).
\]
On fixing the first coordinate, the second is a positive quadratic in
the transverse coordinate. At separation $r\geq\sqrt H$, its derivative
has size comparable to $r$, and uncertainty $O(H)$ permits only
boundedly many interval pairs of width $r$. In the nearest part, the
positive coefficients in the first coordinate already give bounded
overlap. Plancherel at each scale and Cauchy--Schwarz over the
logarithmically many scales reduce Equation~\eqref{eq:pair} to the sum
of the individual spacetime product integrals with the squared time
cutoff.

Localize each initial factor $a_Ig,b_Jg$ further by
$G(x/r-k)$, $k\in\Z$. Here $G$ is Schwartz, has Fourier support in a
sufficiently small neighborhood of zero, and
$\sum_kG(v-k)=1$. Such a function is obtained by choosing its Fourier
transform supported near zero with the normalization prescribed by
Poisson summation. The squared factors have bounded sum. Multiplication
enlarges a velocity band by at most $cH/r\leq cr$, preserving the
separation of separated pairs. Write $e_{I,k}$ for the energy of a
localized factor, including its exact indices. Then
\begin{equation}
 \sum_{I,k}e_{I,k}\lesssim e_{\rm cap},\qquad
 e_{I,k}\lesssim\tau^{-O(\eps)}b_0r^{1-\eta}
                       +\mathrm{negl}\,e.
 \tag{LC}\label{eq:LC}
\end{equation}
For the second bound, use the original boosted input, before cap
selection. The localized operator needs only its analyzed rows with
velocity within $O(\tau^{-\eps}r)$ of $I$ and position within
$O(\tau^{-\eps}r)$ of $rk$. On those rows apply the boundedness of
synthesis, the Fourier multiplier, and the position multiplier, followed
by input regularity. The root packet widths are $\sqrt H\leq r$.
Excluded velocities give Gaussian frequency tails; excluded positions
give Gaussian spatial tails and the Schwartz kernel tail of a multiplier
whose spatial width is $O(H/r)\leq O(r)$. Applying Schur bounds after
rescaling these kernels proves negligible norm for the discarded
transfers. This proves Equation~\eqref{eq:LC} with constants independent
of the exact-index spaces and Hilbert multiplicities.

Truncate the time integral to $|t|\leq\tau^{-\eps}$. The time cutoff's
Schwartz decrease, together with polynomial Bernstein bounds for the
bandlimited fields, makes the error negligible. For a localized piece,
discard the output outside
\[
 |x-rk-tN_I|\leq\tau^{-C\eps}r,
\]
where $C>2$ is fixed large enough. To obtain the $L^2$ error bound, first
discard the position multiplier's tail beyond $\tau^{-\eps}r$, retaining
a smooth propagation plateau equal to one on the original enlarged
Fourier support. The nonstationary kernel tail then applies, since
$r^2\geq H$ and $|t|\leq\tau^{-\eps}$. Tensor-product errors are
negligible by Bernstein on the untruncated factors. These errors sum over
polynomially many significant position indices; the remaining indices
can be discarded together by the original support and tail bounds.
On retained outputs a product with the shifted factor requires
\[
 |k-k'|\lesssim\tau^{-C'\eps},
\]
because $|N_I-N_J|\lesssim r$ and
$|\Delta|\lesssim\tau^{1-\eps}\lesssim\tau^{-\eps}r$.
At each point only $\tau^{-O(\eps)}$ such localized terms occur, so
Cauchy--Schwarz separates them at this loss.

For a separated localized pair, extend its spacetime integral to the
whole plane. Bilinear Plancherel gives
\[
 \int_{\R^2}\norm{u_{I,k}(t,x)}^2
                    \norm{u_{J,k'}(t,x+\Delta)}^2\,dx\,dt
 \lesssim r^{-1} e_{I,k}e_{J,k'}.
\]
Indeed the map from momenta to their summed spatial and temporal
frequencies has multiplicity at most two and Jacobian
$|Hp_1/\sigma-Hp_2/\sigma'|=|N_1-N_2|\gtrsim r$.
The shift introduces only a phase. This calculation is componentwise
valid for Hilbert-valued fields and then integrates over the independent
exact indices. The shift interval cancels the prefactor $\tau^{-1}$
in Equation~\eqref{eq:pair}, up to $\tau^{-O(\eps)}$.

There are two possibilities for the nearest scale. If $r_0=\sqrt H$,
spatial Bernstein has bandwidth $O(r_0/H)=O(r_0^{-1})$; combine it with
the other factor's conserved $L^2$ norm over the retained times. If
$r_0=\tau$, extend the shift integral to all of $\R$, after imposing the
position-pair restriction just proved. The spatial and shift integrals
then give the product of the two energies, and the surviving prefactor
is $\tau^{-1}=r_0^{-1}$. Thus all cases are bounded by
\[
 \tau^{-O(\eps)}r^{-1}e_{I,k}e_{J,k'}.
\]
Each localized piece has only $\tau^{-O(\eps)}$ allowable partners.
Use Equation~\eqref{eq:LC} for one energy and sum the other to obtain
$\tau^{-O(\eps)}r^{-\eta}b_0e_{\rm cap}$ at each scale. Since
$r\geq\tau$, summing scales and caps proves the free estimate after
letting $\eps$ tend to zero. In particular this argument does not replace
a finite-time spatial mass by an unrestricted spacetime $L^4$ norm when
$H$ is very small.

It remains to insert masks. Induct on their number. At the first cut of
length $\ell$ carrying a mask, peel the raw analyzed rows into dyadic
scalar-regularity levels. At threshold $b$, use the tests with widths
$f,\ell f$ to remove disjoint heavy parts of mass greater than
$bf^{1-\eta}$. The preceding threshold ensures that their union has
regularity $O(b)$. Include an arbitrarily low polynomial floor remainder.
If $M$ is the row contraction, write it as the coherent sum
$\sum_bM\mathbf1_{\mathrm{level}\ b}$ before synthesis. The free
estimate to this cut shows that its level's summed raw mass $z_b$
satisfies
\[
 b z_b\lesssim s^{-o(1)}\ell^{-\eta}b_0e.
\]
The post-mask energies sum to at most $z_b$, and same-frame Gaussian
locality gives regularity $O(b)$ after reanalysis. Apply the induction
hypothesis separately on each descendant experiment, whose relative
terminal time length is $\tau/\ell$, and sum its energies. The scale
factors multiply as
$\ell^{-\eta}(\tau/\ell)^{-\eta}=\tau^{-\eta}$.
The floor has an arbitrarily small regularity constant; a trivial total
energy bound makes its contribution negligible. Post-states too small to
satisfy polynomial bounds relative to their own energy are instead
estimated directly on the original polynomial terminal tests.

There are only logarithmically many levels. For fixed assignment
multipliers the fourth root of the weighted mass-square sum is an
$\ell^4$ norm of Hilbert norms and hence a seminorm in the input state.
The triangle inequality therefore sums the coherent levels at a
subpower loss. Repeating this a fixed number of times completes the
induction. Repeated cuts and subpower time gaps use the same-frame bound
and the elementary estimate already proved, with subpower loss.
\end{proof}

\subsection{Memory of the base angle}

\begin{lemma}[Memory bound for documented labels]\label{low:memory}
Use an affine documented configuration as in Equation~\eqref{eq:F},
with ordinary regularity imposed immediately before the relevant gates.
Fix $0<i=j-h<j\leq1$, set $\delta=s^h$, and, separately for each starting
interval, normalize time length to one and the ending widths to
$r_j=f_j=1$. In the quadratic case set $Q=q$ and $C=q-J$; in the linear
case set $Q=0$ and $C=c_1$. Suppose
\begin{equation}
 L\geq1,\qquad 2L\geq1+\max(1,Q),\qquad L\geq1+C.
 \tag{ML}\label{eq:ML}
\end{equation}
Assign to each $j$-label a list of at most $\delta^{-R}$ cells of width
$\delta^L$ for the normalized base angle $X$. Then
\begin{equation}
 \Prob^j\{X\text{ belongs to its label's list}\}
 \leq s^{-o(1)}
       \delta^{(1-d+2\Gamma+\eta(L-1)-R)/2}.
 \tag{mem}\label{eq:mem}
\end{equation}
The cell origins may depend on the label. The assertion of power rarity
requires a strictly positive exponent after all outer errors.

In the classical model there is also a prediction version. Drop
$L\geq1+C$ from Equation~\eqref{eq:ML}. For each exact ancestor row at
$i$, allow at most $s^{-o(1)}\delta^{-Z}$ candidate values for the ending
derivative $P_j(F)$ on its true base ray, in the normalized units. Restrict
the event to guesses accurate to $s^{-o(1)}\delta^{1-L}$. Then
Equation~\eqref{eq:mem} holds with an additional $-Z$ in its exponent's
numerator. In the linear case one may guess the enclosing test's linear
coefficient instead.
\end{lemma}

\begin{proof}
Partition the exact base variables at the initial time into lattice
boxes $M$ of widths $(\delta^L,\delta^L,\delta^{2L})$, using the tilt
of the center of each $X$-bin for the last coordinate. These variables
do not mix under propagation. For a fixed listed angle cell, a terminal
label permits initial horizontal and tilted vertical positions in
intervals of width $O(\delta)$ each. To check this, backflow its terminal
base footprint through a bounded time; variation within the angle bin
contributes $O(\delta^L)$ horizontally and $O(\delta^{2L})$ vertically
in the bin's tilt, while the terminal footprint has width $O(\delta)$
in both relevant directions. A change from the label's tilt to the
bin's tilt has bounded coefficients. Therefore at most
\[
 s^{-o(1)}\delta^{2-3L-R}
\]
base boxes $M$ are relevant to any ending label and its whole list.

On the initial gate within $M$, bin both curvature coefficients, at its
starting time and in segment units, to width
$\delta^{-\max(1,Q)}$. Equation~\eqref{eq:A} shows that each ending
label knows only $s^{-o(1)}$ possible bins, including the extrapolation
of its quadratic coefficient back to the initial time. Subtract the
shear of that rounded curvature trajectory. In the linear case omit
this operation and use the enclosing linear tests supplied by the
affine reduction. Residual curvatures of the two endpoint tests are now
bounded by $s^{-o(1)}\delta^{-\max(1,Q)}$ throughout the segment.

The endpoint derivatives at a true ray differ by at most
$s^{-o(1)}\delta^{-\max(0,C)}$. Recenter at the ray $F_M$ through the
base-box center; this adds at most
$s^{-o(1)}\delta^{L-\max(1,Q)}$. Both quantities are bounded by
$s^{-o(1)}\delta^{1-L}$, by Equation~\eqref{eq:ML}. Bin the residual
initial derivative at that accuracy and subtract the resulting linear
shear. The ending label, $M$, and the curvature bins determine only
$s^{-o(1)}$ possibilities for this last bin. The residual ending
derivative at $F_M$ is consequently at most
$s^{-o(1)}\delta^{1-L}$.

For the prediction version, clip the lists to a polynomial parameter
range large enough to contain all accurate guesses. Subtract the
rounded curvature's derivative at the true ray from each guess, and bin
the result to width $\delta^{1-L}$. Make one copy of the classical
ancestor row for each distinct resulting bin. This costs at most
$s^{-o(1)}\delta^{-Z}$ in total energy. On the event of an accurate
guess, recentering at $F_M$ has the same error as before. The ending
label and $M$ again know only $s^{-o(1)}$ pertinent bins. If the guess
uses the true derivative but the test is enclosed by a linear one,
their coefficients differ by at most $s^{-o(1)}$, which fits the
tolerance. Copying is used only in the classical model, where energy
contributions add pointwise over exact indices.

After these subtractions, the part of each ending label over $M$ lies
in an ordinary scalar test of velocity width $s^{-o(1)}$ and position
width $\delta s^{-o(1)}$. To verify the tighter position statement,
Taylor-expand the two shear polynomials about $F_M$. Their horizontal
linear terms are bounded by
$\delta^{1-L}\delta^L=\delta$. The remaining terms are bounded by
$\delta^{-\max(1,Q)}\delta^{2L}\leq\delta$.
The tilted vertical coordinate has precisely the latter order, also
after base free evolution. The velocity estimate follows from the
same expansion with its less restrictive unit tolerance.

For each starting time and bin tuple, ordinary regularity before the
gate gives scalar input regularity constant
\[
 s^{-o(1)}\kappa_i w_j\delta^{(3+\eta)L}.
\]
Indeed a scalar regularity box pulled back by the subtracted common
shear, together with the base box $M$, is a full cylinder test of base
width $\delta^L$. Restriction to the gate and same-frame reanalysis
preserve this upper bound. Apply Proposition~\ref{low:scalar} to each
such experiment, allowing all subsequent masks and retaining exact
indices. The energies sum to at most $s^{-o(1)}E_i$, or to
$s^{-o(1)}E_i\delta^{-Z}$ in the prediction version. The scalar cutoff
conditions follow from the original density and phase cutoff bounds.
As throughout, extremely small-energy pieces are bounded directly on
the polynomial output tests.

Let $m_{\lambda,M}$ be the actual read mass of the tested event.
Each pair $(\lambda,M)$ sees only $s^{-o(1)}$ bin tuples. Other tuples
have no allowed packet chain and contribute negligible operator norm.
Cauchy--Schwarz in the contributing fields and
Equation~\eqref{eq:SC} imply
\[
 \sum_{\lambda,M}m_{\lambda,M}^2
 \leq s^{-o(1)}\kappa_iw_jE_i
                    \delta^{(3+\eta)L-\eta-Z}.
\]
There is no coherent loss over $M$, which is an exact base partition.
In the classical prediction version use only accurate copied parts;
they majorize the tested event pointwise and have the same small-list
property.

The documented count and the base-box count give at most
\[
 s^{-o(1)}\frac{E_j}{\kappa_jw_j}
                        \delta^{2-3L-R}
\]
pairs. Apply Cauchy--Schwarz to their masses and divide by $E_j^2$.
Since $E_j/E_i=\delta^{-d+o(1)}$ and
$\kappa_j/\kappa_i=\delta^{p+o(1)}$, the exponent is
$d-p+2+\eta(L-1)-R-Z$. Finally
$d-p+2=1-d+2\Gamma$, giving the assertion.
\end{proof}

\begin{corollary}[Memory when both costs vanish]\label{low:zero-memory}
If $c_2=c_1=0$, Equation~\eqref{eq:mem} holds with $j=1$ and $L=1$
for every fixed terminal depth gap, without assuming affine profiles.
\end{corollary}

\begin{proof}
In normalized tail units the leaf curvature is at most
$s^{-o(1)}\delta^{-1}$ throughout the segment. Its derivative on its
own base footprint is at most $s^{-o(1)}$, by the two zero-cost bounds
at the leaf. Recentring on a compatible base box of width $\delta$
changes the derivative by only $s^{-o(1)}$. Thus the Taylor bounds in
the preceding proof hold directly, without curvature or derivative
bins. The scalar-channel and counting arguments are unchanged.
\end{proof}

\subsection{The zero-cost hybrid case}

\begin{proposition}[Exclusion of a strictly nonclassical zero-cost case]
\label{low:hybrid-zero}
Suppose $\Gamma_{\rm hyb}>0$ and
$\Gamma_{\rm hyb}>\Gamma_{\rm cl}$. The minimizing hybrid families
cannot have $c_2=c_1=0$.
\end{proposition}

\begin{proof}
Write $\Gamma=\Gamma_{\rm hyb}$ and assume both costs vanish. Normalize
the leaf widths to one. The bottom Planck assertion
\eqref{eq:Pl} gives $H\leq s$ and $H=s^{1+o(1)}$. At each leaf,
\[
 |\beta(T)|+s|K|\leq s^{-o(1)},\qquad
 |P(x,b)|\leq s^{-o(1)}.
\]
Taylor expansion on its base footprint, whose horizontal and tilted
vertical position widths are $s$, shows that its normal velocity and
position lie in unsheared intervals of widths $s^{-o(1)}$ and
$s^{1-o(1)}$ respectively.

\paragraph{Initial localization.}
Partition $X$ into unit bins. Split the initial normal field by a smooth
compactly supported unit Fourier-velocity partition of unity, and then
by unit-width Schwartz position factors forming a partition of unity
and having bounded Fourier support. Summed energies are $O(e)$.
Every leaf sees only $s^{-o(1)}$ groups: its angle and normal velocity
windows have subpower widths, and its position backflows into the
corresponding subpower neighborhood of an initial position bin.
This follows through the fixed stack of packet kernels since
$d_s\leq1$. Gaussian and Schwartz tails justify deleting all other
transfers with negligible norm; first truncate to polynomially many
significant input bins. Equivalently use arbitrary fixed-power
enlargements and then send their exponents to zero.

By the triangle inequality on each assignment, it suffices to estimate
the groups separately and sum their cube sums at subpower loss.
Tiny-energy groups admit direct polynomial-test bounds. On a remaining
group boost and translate so that $|X|\lesssim1$, the normal Fourier
velocities are bounded, and the initial normal position has negligible
mass outside $|D|\leq s^{-o(1)}$. Denote its energy by $e_g$.

The group has input regularity $O(\kappa)$. The same assertion holds
after replacing it by the direct sum of a smooth Fourier-velocity
partition at any width $r\gg\sqrt H$. Here is the operator justification
needed later. The conjugated operators in the root packet frame have
rapidly decreasing Schur kernels in its normal packet units. For a
vector of Fourier multipliers the $\ell^2$ norm of its derivative of
order $k$ is $O(r^{-k})$. Rescale the Gaussian frequency-side inner
product to root packet units and integrate by parts; since
$r\gg\sqrt H$, this gives a bound for the whole vector kernel with
summable tails. Unit position factors have the analogous space-side
bound. At each exact base point these operators need only bounded
enlargements, and then summable enlarged shells, of a normal test.
Same-frame reanalysis therefore proves the claimed regularity for the
sum measure. Integrating over the subpower normal support gives the base
energy marginal regularity
\[
 \mu_{\rm base}(\text{base test of width }R)
 \leq s^{-o(1)}\kappa R^3,
 \qquad s^2\lesssim R\lesssim1.
\]
We used $R^{3+\eta}\leq R^3$ and covered normal support by subpower
many unit windows. All these weights are polynomial, so tails remain
negligible.

After the boost, base leaf positions lie in squares of side
$s^{1-o(1)}$. Pass from assignments to lattice position masses at
subpower loss and drop angular restrictions. Let $y$ denote all exact
indices, $Y_d(T)$ the indices whose base position at $T$ belongs to a
square $d$ of side $s$, and $u_y(T,x')$ the output field before final
analysis. The unnormalized functional is bounded, up to subpower loss,
by
\begin{equation}
 \sum_{T,d}\int_\R
 \left(\int_{Y_d(T)}\norm{u_y(T,x')}^2\,dy\right)^{3/2}dx'.
 \tag{HL}\label{eq:HL}
\end{equation}
Indeed the analyzer's spatial marginal is convolution by a probability
Gaussian. Jensen removes that convolution in this upper bound, and
H\"older on a normal bin of length $s$ cancels the functional's factor
$s^{-1/2}$. We may replace the fields by approximations having negligible
$L^2$ error at every time: the original assignment functional absorbs
these errors by its polynomial parameters. We may likewise restrict
the output to a polynomial range of $x'$ when necessary.

\paragraph{A stopping scale and its envelopes.}
Fix $0<\zeta<1/4$, set $r\asymp s^\zeta$ dyadically, and choose
comparable dyadic scales
\[
 \ell_0=H/r^2,\qquad B_*=H/r.
\]
Then $1\gg\ell_0\gg B_*\gg s$ on sufficiently accurate outer
approximations, and
$d_{\ell_0}\asymp r$, $q_{\ell_0}\asymp B_*$.
We will bound the expression in Equation~\eqref{eq:HL} by
\[
 s^{-o(1)}\bigl(r^{-D_0}+B_*^{-\Gamma}\bigr)e_g\sqrt\kappa,
\]
where $D_0$ is absolute and independent of the number of masks.
Separated velocity caps give the first term: a transverse crossing
estimate will be combined with the kinetic bound. Nearby caps give the
second term: we compare them with one experiment read at time length
$B_*$, apply the exponent $\Gamma$ to its prefixes, and use the kinetic
bound on each remaining gap. Since $B_*=s^{1-\zeta+o(1)}$, both terms
improve on the extremal exponent once $D_0\zeta<\Gamma$.

Split the group input coherently into $O(r^{-1})$ smooth velocity caps
of width $r$ centered at bounded lattice points $\theta_k$. Call their
outputs $u_{k,y}$. Let $e_y$ be the original group fiber energy, so
$\int e_y\,dy=e_g$ and each cap fiber has energy $O(e_y)$.

Add the cut $\ell_0$, including any masks already at that length. For a
fixed stack of $n$ stages use slack $s^{-\eps}$, with $\eps>0$
arbitrarily small. At each mask restrict its rows to velocities within
$C_ns^{-\eps}\max(r,d_\ell)$ of the cap center, and post-compose with
a smooth projection onto a slightly larger Fourier band. Choose
successively increasing constants to include the preceding supports.
The modified operators have bounded norm, and Gaussian frequency tails
make their output differences negligible in $L^2$, uniformly per
fiber and branch. At the stopping interval $I$, write their states as
$w_{I,k,y}$; their bandwidths in velocity are $O_n(s^{-\eps}r)$.
The fixed factors $s^{-C_n\eps}$ below become subpower losses by taking
the inner limit before $\eps\downarrow0$, even if the mask counts grow
in a subsequent outer approximation.

For $T\in I$ the modified fields satisfy
\begin{equation}
 \norm{u_{k,y}(T,x')}^2
 \lesssim_n\frac{s^{-C_n\eps}}{B_*}
 \int_{|x'-z-(T-T_I)\theta_k|\leq s^{-C_n\eps}B_*}
             \norm{w_{I,k,y}(z)}^2\,dz
 +\mathrm{negl}\,e_y.
 \tag{env}\label{eq:env}
\end{equation}
To prove this, insert a smooth plateau on the stopping state's band
before its first propagation. For duration at most $\ell_0$, its
kernel has size $O(s^{-C_n\eps}/B_*)$ and rapid tails beyond the
displayed displacement. After the stop, $d_\ell\gtrsim r$ and
$q_\ell\lesssim B_*$. The position kernel of a mask restricted to row
velocities of width $O_n(s^{-\eps}d_\ell)$ has absolute operator-valued
bound $O_n(s^{-\eps}/q_\ell)$ and Gaussian tails on scale $q_\ell$:
insert the two normalized analyzing packets, bound the row contraction
by one, and integrate their product against the frame measure.
Post-projection and bandlimited free evolution for a duration at most
$\ell$ have Schwartz tails on the same scale, with slack. The effective
uncertainty parameter is bounded in $(q_\ell,d_\ell)$ units. Each
kernel has integral bounded by $s^{-C_n\eps}$. Compose after the first
convolution and apply Cauchy--Schwarz against its absolute kernel.
The resulting rapid tails can be bounded by negligible times $e_y$.
This proves Equation~\eqref{eq:env}, also for repeated cuts.

Expand the squared norm of $\sum_ku_{k,y}$ in
Equation~\eqref{eq:HL}. Pairs with centers within $s^{-2\eps}r$
are bounded by $s^{-O(\eps)}\sum_k\norm{u_{k,y}}^2$.
Call the remaining pairs high pairs, and write
$V=|\theta_k-\theta_{k'}|\geq s^{-2\eps}r$ for their separation.
Only their sum will incur an ordinary power of $r^{-1}$.

\paragraph{The transverse crossing estimate.}
For normal fields $f,g$ and $a>0$, define the crossing energy
\[
 \mathcal C_a(f,g)=\int_{|x_1-x_2|\leq a}
                       \norm{f(x_1)}^2\norm{g(x_2)}^2\,dx_1dx_2.
\]
For each high pair and exact index,
\begin{equation}
 s\sum_T\int_\R
       \norm{u_{k,y}(T,x')}^2\norm{u_{k',y}(T,x')}^2\,dx'
 \lesssim s^{-O_n(\eps)}V^{-1}e_y^2+\mathrm{negl}\,e_y^2.
 \tag{cross}\label{eq:cross}
\end{equation}
We prove the localization needed before the stopping cut, so that
independent estimates on many short intervals do not introduce a scale
loss. At each cut at or before the stop call the two modified band
states $g_{1,I},g_{2,I}$, and put $W_I=\ell_I V$. For sufficiently large
fixed $D$, their stopping crossing energy satisfies
\[
 \sum_I\int_{|x_1-x_2|\leq s^{-D\eps}W_I}
       \norm{g_{1,I}(x_1)}^2\norm{g_{2,I}(x_2)}^2\,dx_1dx_2
 \lesssim s^{-O_{D,n}(\eps)}e_y^2.
\]
We establish this by pulling back successively from each cut to its
parent, increasing the slack constant a fixed amount at each step.
At the root total energy gives the assertion.

Work in a common Galilean frame where both band velocities are $O(V)$.
They remain separated by a constant times $V$ before the stop.
For a parent $I$ and its children $J$ at the next cut write
$W=\ell_IV$, $w=\ell_JV$, and
\[
 g_{\alpha,J}=A_{\alpha,J}U(T_J-T_I)g_{\alpha,I}.
\]
The required one-step estimate is
\[
 \sum_{J\subset I}\mathcal C_{s^{-D\eps}w}(g_{1,J},g_{2,J})
 \leq s^{-O_{D,n}(\eps)}
       \mathcal C_{s^{-O_{D,n}(\eps)}W}(g_{1,I},g_{2,I})
       +\mathrm{negl}\,\norm{g_{1,I}}_2^2\norm{g_{2,I}}_2^2.
\]
Thus the child crossings must be charged to nearby pairs in the parent,
with no factor for the number of children. We first control the position
spread of the cut operators, then sum the freely evolving separated
pairs by bilinear Plancherel.

Decompose each parent state into
\[
 g_{\alpha,I}
 =\sum_{b\in\Z}P_\alpha
                  (\mathbf1_{[bW,(b+1)W)}g_{\alpha,I}),
\]
where $P_\alpha$ is a smooth plateau of width comparable to $V$,
equal to one on the input band and with the two enlarged supports still
separated. At a child crossing, only pairs with
$|b-b'|\leq s^{-O_D(\eps)}$ contribute, up to negligible error.
Indeed free band propagation has negligible tail beyond a slack
multiple of $W$, and the child masks and projections have position
width at most a slack multiple of $w$. The needed inequalities are
\[
 H/V\lesssim\ell_JV,
 \qquad q_{\ell_J}\lesssim\ell_JV,
 \qquad H/r\lesssim\ell_JV,
\]
which follow from $\ell_J\geq\ell_0$ and $V\gg r$.
Here is a quantitative tile argument for this step. For intervals
$E_m=[mw,(m+1)w)$, put
\[
 B_{\alpha,J,t}=A_{\alpha,J}U(T_J-t)\widetilde P_\alpha,
 \qquad
 K_{mn}=\norm{\mathbf1_{E_m}B_{\alpha,J,t}\mathbf1_{E_n}}_{2\to2},
\]
where $\widetilde P_\alpha$ is a slightly larger plateau equal to one
on the support of $P_\alpha$. Its two supports remain separated by
a constant times $V$. Uniformly in $J$, $t\in J$, and the measurable
row contractions, $\norm{B_{\alpha,J,t}}_{2\to2}\leq C_n$, and there
is $S_0=s^{-C_n\eps}$ such that, for every $N$ and $R\geq 2S_0$,
\begin{equation}
 \sup_m\sum_{|n-m|>R}K_{mn}
 +\sup_n\sum_{|m-n|>R}K_{mn}
 \leq C_{N,n}S_0^{C_n}(R/S_0)^{-N}.
 \label{low:position-matrix}
\end{equation}
We verify that this bound requires no derivatives of a mask. For the
length-$\ell$ frame, Plancherel in packet velocity gives, for any
position set $E$ and $L>0$,
\[
 \norm{\mathbf1_{\{D:\dist(D,E)>Lq_\ell\}}
                    V_\ell\mathbf1_E}_{2\to2}
 \leq C e^{-cL^2}.
\]
Indeed its squared norm is bounded by the supremum, for $x\in E$, of
the Gaussian integral over those packet centers $D$. To bound a mask
between separated position sets, split packet centers into those
within one third of their separation from the input set and the
complement. On the complement use this analysis bound; on the first
part use its adjoint for synthesis into the output set. The row
contraction preserves both sets of centers. Thus the off-diagonal
mask norm has Gaussian decrease on scale $q_\ell$, uniformly in the
mask and Hilbert multiplicity.

The Fourier projections have Schwartz kernel tails on scale $H/r$.
For $U(T_J-t)\widetilde P_\alpha$, velocities are $O(V)$ and travel
at most $O(w)$, while the kernel's wavelength is $O(H/V)\leq O(w)$.
Outside a sufficiently large multiple of $w$, its phase derivative
is bounded below by a constant times position displacement.
Integration by parts gives arbitrarily rapid integral tails in
displacement divided by $w$. All three widths are at most $w$ by
the displayed inequalities. Products of the finite number of such
operators have the matrix tail bound \eqref{low:position-matrix}:
at large total tile separation one factor has comparably large
separation, and convolution of their summable tails preserves
arbitrary decay. The slack enlargements only replace $w$ by $S_0w$
and insert a fixed power of $S_0$.

Choose $R=s^{-C'_{D,n}\eps}$ with $C'_{D,n}$ larger than the constant
in $S_0$. Equation~\eqref{low:position-matrix}, with $N$ arbitrarily
large after fixing $\eps$, allows deletion of all matrix blocks with
$|m-n|>R$ in negligible operator norm. If $B^R$ is the retained
operator, Cauchy--Schwarz and the individual block bound $K_{mn}\leq C_n$
give the explicit local inequality
\[
 \norm{\mathbf1_{E_m}B^Rf}_2^2
 \leq C_n(2R+1)\sum_{|n-m|\leq R}
                              \norm{\mathbf1_{E_n}f}_2^2.
\]
For $S\geq1$, cover the crossing region by tile pairs
$|m-m'|\leq S+2$ and apply the local inequality to both factors.
Each input pair occurs at most $(2R+1)^2$ times. The triangle
inequality between tile indices therefore gives
\begin{equation}
 \mathcal C_{Sw}(B_1f_1,B_2f_2)
 \leq C_n(2R+1)^4
        \mathcal C_{(S+2R+5)w}(f_1,f_2)
       +\mathrm{negl}\,\norm{f_1}_2^2\norm{f_2}_2^2.
 \label{low:product-transfer}
\end{equation}
The error estimate follows by viewing $\mathcal C_a^{1/2}$ as the
$L^2$ norm of a restricted tensor product and using the negligible
operator-norm difference between $B$ and $B^R$.

Apply the same matrix argument at tile width $W$ to
$A_{\alpha,J}U(T_J-T_I)P_\alpha$. All its movement and kernel widths
are at most a slack multiple of $W$. After negligible deletion, on
each output $W$-tile there are only $s^{-O_{D,n}(\eps)}$ contributing
parent pieces. Pointwise Cauchy--Schwarz for their sums in both
factors bounds the child crossing by the sum of the crossings of
individual parent piece pairs, with a slack factor. A contributing
pair must have $|b-b'|\leq s^{-O_{D,n}(\eps)}$: the output positions
are within $s^{-D\eps}w\leq s^{-D\eps}W$ and both input tile indices
are within a slack distance of their output tile. Dropping the output
tile restrictions on each retained pair only increases its crossing
integral. This proves the earlier parent-pair restriction with its
required restricted-region bookkeeping.

For a retained pair write
$f_{\alpha,b}=P_\alpha(\mathbf1_{[bW,(b+1)W)}g_{\alpha,I})$ and
$v_{\alpha,b}(t)=U(t-T_I)f_{\alpha,b}$. At every $t\in J$,
$B_{\alpha,J,t}v_{\alpha,b}(t)$ is exactly the cut image of this
piece. Use Equation~\eqref{low:product-transfer} with
$S=s^{-D\eps}$, and write $L=s^{-O_{D,n}(\eps)}$ for its enlarged
radius and prefactor. All next-cut children have the same length
$\ell_J$ and are disjoint. Hence averaging over each child gives
\begin{align}
 &\sum_{J\subset I}\mathcal C_{Sw}
       \bigl(A_{1,J}v_{1,b}(T_J),A_{2,J}v_{2,b'}(T_J)\bigr)
 \notag\\
 &\qquad\leq \frac{L}{\ell_J}
    \int_I\int_{|\Delta|\leq Lw}\int_\R
       \norm{v_{1,b}(t,x)}^2\norm{v_{2,b'}(t,x+\Delta)}^2
                  \,dx\,d\Delta\,dt
       +\mathrm{negl}\,\norm{f_{1,b}}_2^2\norm{f_{2,b'}}_2^2
 \notag\\
 &\qquad\leq C\frac{L^2w}{\ell_JV}
                  \norm{f_{1,b}}_2^2\norm{f_{2,b'}}_2^2
       +\mathrm{negl}\,\norm{f_{1,b}}_2^2\norm{f_{2,b'}}_2^2.
 \label{low:child-crossing-sum}
\end{align}
The last line extends the time integral to $\R$ and uses bilinear
Plancherel for the separated plateau supports, independently for each
shift. Thus the only scale factor is $w/(\ell_JV)=1$; there is no
factor equal to the number of children.

Finally, $\norm{f_{\alpha,b}}_2^2\leq
C\norm{\mathbf1_{[bW,(b+1)W)}g_{\alpha,I}}_2^2$.
If $|b-b'|\leq K_0=s^{-O_{D,n}(\eps)}$, every pair in these original
tiles satisfies $|x_1-x_2|\leq(K_0+2)W$. The tiles are disjoint, so
\[
 \sum_{|b-b'|\leq K_0}
       \norm{f_{1,b}}_2^2\norm{f_{2,b'}}_2^2
 \leq C\mathcal C_{(K_0+2)W}(g_{1,I},g_{2,I}).
\]
Combining this with Equation~\eqref{low:child-crossing-sum} proves
the one-step pullback into the enlarged, restricted parent crossing
energy. Finite induction proves the stopping estimate. Polynomially
many intervals multiply only negligible errors, whose decay order was
chosen after all cutoff powers were fixed.

\paragraph{From stopping energy to terminal crossings.}
Now insert Equation~\eqref{eq:env} for both outputs from a stopping
interval. For fixed $z_1,z_2$, their kernel cores overlap in a position
interval of length at most $s^{-O_n(\eps)}B_*$, and at only
$s^{-O_n(\eps)}B_*/(sV)$ terminal times. We used $B_*\gg s$ to absorb
the possible extra endpoint time. They overlap only if
$|z_1-z_2|\leq s^{-O_n(\eps)}\ell_0V$. The two envelope prefactors,
the position length, and $s$ times the time count give
$s^{-O_n(\eps)}V^{-1}$. The stopping crossing bound proves
Equation~\eqref{eq:cross}. A constant additive tail from one envelope
is integrated using the other's $L^2$ bound, or the spatial integral of
its core, so unbounded position causes no extra error.

\paragraph{Summing the high pairs.}
At a given time put
$H_y(x')=\norm{u_{k,y}(T,x')}\norm{u_{k',y}(T,x')}$, and define
\[
 M_{T,d}=\int_{Y_d(T)}e_y\,dy,\qquad
 N_{T,d}=\int_{Y_d(T)}e_y^{-1}
                    \int_\R H_y(x')^2\,dx'\,dy,
\]
with zero integrand when $e_y=0$. Contraction gives
$\int\int_{Y_d(T)}H_y\,dy\,dx'\lesssim M_{T,d}$, while
Cauchy--Schwarz in $y$ bounds its squared $L^2$ norm by
$M_{T,d}N_{T,d}$. Interpolation, or Cauchy--Schwarz between these two
integrals, therefore gives
\[
 \int_\R\left(\int_{Y_d(T)}H_y(x')\,dy\right)^{3/2}dx'
 \lesssim M_{T,d}N_{T,d}^{1/2}.
\]
Lemma~\ref{low:kinetic} applied to the original group's base marginal
gives $\sum M_{T,d}^2\leq s^{1-o(1)}\kappa e_g$.
Equation~\eqref{eq:cross} gives
$\sum N_{T,d}\lesssim s^{-1-O_n(\eps)}V^{-1}e_g$.
Cauchy--Schwarz in $(T,d)$ bounds one high pair by
$s^{-o(1)-O_n(\eps)}V^{-1/2}e_g\sqrt\kappa$.
There are $O(r^{-2})$ pairs, and the finite-sum inequality for the
$3/2$ power introduces only another fixed power of their number.
Hence all high terms in Equation~\eqref{eq:HL} are bounded by
\[
 s^{-o(1)-O_n(\eps)}r^{-D_0}e_g\sqrt\kappa
\]
for an absolute $D_0$, independent of the number of masks.

\paragraph{Summing the low contribution.}
Partition each stopping interval into intervals $J_*$ of length $B_*$,
and let $L_b=[bB_*,(b+1)B_*)$. Use one actual cylinder experiment whose
root is the direct sum of the $r$-cap inputs, whose masks through the
stopping cut act componentwise, and which thereafter evolves freely.
For each $J_*$, let $\nu_{J_*}$ be this experiment's analyzed row measure
at its left endpoint. All $J_*$ and all position bins $L_b$ thus belong
to the same experiment; its root energy is $O(e_g)$ once, not the sum
of copied energies over $J_*$. The analytical
approximation operators used to prove the envelopes are not included
in this experiment. For $T\in J_*$ and $x'\in L_b$,
\[
\begin{split}
 \int_{Y_d(T)}\sum_k\norm{u_{k,y}(T,x')}^2\,dy
 \leq{}&\frac{s^{-O_n(\eps)}}{B_*}
 \nu_{J_*}\bigl\{y\in Y_d(T),\ |N|\leq C,\\
 &\hspace{38mm}\dist(D_{T_{J_*}},L_b)
                  \leq s^{-O_n(\eps)}B_*\bigr\}
 +\mathrm{negl}\,e_g.
\end{split}
\]
To prove this, move the windows in Equation~\eqref{eq:env} from the
stopping time to $T_{J_*}$. The inverse bandlimited transfer needs
displacement uncertainty at most
$\ell_0O_n(s^{-\eps}r)=O_n(s^{-\eps}B_*)$, so local stopping energy
is bounded by the freely continued energy on the displayed enlarged
window. Pass to analyzed mass using its spatial Gaussian marginal,
since $q_{B_*}\lesssim B_*$. Gaussian frequency tails allow restriction
to $|N|\leq C$, because $d_{B_*}\asymp\sqrt r$ and cap centers are
bounded. Take $\eps$ small relative to $\zeta$ and the fixed number
of cuts. The modified stopping states differ negligibly from the
actual direct-sum states. Cover enlarged windows by
$s^{-O_n(\eps)}$ ordinary $B_*$-bins, with the same overlap bound.

Integrate in $x'$; a polynomial restriction is allowed as explained
after Equation~\eqref{eq:HL}. The low contribution is thus, up to slack
and negligible errors, at most
\begin{equation}
 B_*^{-1/2}\sum_{J_*,b}\sum_{T\subset J_*,d}
 \left[\nu_{J_*}\{D_{T_{J_*}}\in L_b,\ |N|\leq C,
                                      \ y\in Y_d(T)\}\right]^{3/2}.
 \tag{low}\label{eq:low}
\end{equation}
Peel ordinary cylinder regularity levels of the positive measure
$\nu_{J_*}$, and let $z_\beta$ be summed mass at threshold $\beta$.
Within each level and each interval the regularity is $O(\beta)$.
The heavy pieces are disjoint prefix assignments across the full
$J_*$ family, so the definition of $\Gamma$ applied once to this
actual direct-sum experiment gives
\[
 \sum_\beta z_\beta\sqrt\beta
 \lesssim s^{-o(1)}B_*^{1/2-\Gamma}e_g\sqrt\kappa.
\]
Its input energy is $O(e_g)$ and its input regularity is $O(\kappa)$,
by the direct-sum bound proved above. The scale $B_*$ has fixed positive
depth and the cutoff powers and number of masks are fixed first.
An arbitrarily low floor contributes negligibly by total mass.

Split Equation~\eqref{eq:low} by these logarithmically many levels.
For fixed $J_*,b,\beta$, restrict to its position and velocity window
and take the base marginal. In units of duration $B_*$ it has
regularity at most $C\beta R^{3+\eta}\leq C\beta R^3$ for $R\leq1$:
apply full cylinder regularity with normal width $O(1)$. This width is
admissible because $d_{B_*}\asymp\sqrt r$. Apply
Lemma~\ref{low:kinetic} with $\delta\asymp s/B_*$. The sum of squares
is at most $s^{-o(1)}\delta\beta$ times the marginal's mass, and the
sum of masses is at most $O(\delta^{-1})$ times that mass.
Cauchy--Schwarz bounds the corresponding $3/2$ sum by
$s^{-o(1)}\sqrt\beta$ times the marginal's mass. Summing the bins and
levels and using the preceding prefix bound gives
\[
 s^{-o(1)-O_n(\eps)}B_*^{-\Gamma}e_g\sqrt\kappa
\]
for the low term. No ordinary power of the number of caps was used here.

Finally sum the original groups and let the inner slack losses vanish.
Since $B_*=s^{1-\zeta+o(1)}$, the low term has exponent strictly below
$\Gamma$. Choose $\zeta>0$ so small that $D_0\zeta<\Gamma$; the high
term also has exponent strictly below $\Gamma$. The outer zero-cost
and Planck defects only contribute $s^{-o(1)}$. Their use in passing
from labels to grids and groups has bounded multiplicity independent
of the number of cuts; choosing the stopping scale from the actual $H$
avoids accumulating its defect along the masks. The resulting strict
improvement contradicts extremality and proves the proposition.
\end{proof}

\subsection{Planar incidences and classical labels of zero cost}
\label{low:planar}

We next exclude simultaneous zero quadratic and linear cost in the classical
cylinder model. We first formulate the planar incidence input so that its
application to conditional probability measures is explicit.

\begin{theorem}[Discretized planar Furstenberg estimate]
\label{low:ren-wang}
Let $0<a\leq 1$, $0<b\leq 2$, and $\omega>0$. There are
$\sigma>0$ and $\Delta_0>0$ with the following property. Work in fixed bounded
coordinate charts for points and affine lines in the plane. Suppose that a
probability measure $\nu$ on lines satisfies
\[
 \nu(B(\ell,r))\leq \Delta^{-\sigma}r^b
 \qquad (\Delta\leq r\leq 1).
\]
For every line in its support, suppose that a probability measure $\mu_\ell$
on a set $P_\ell$ within distance $C\Delta$ of that line satisfies,
uniformly in $\ell$,
\[
 \mu_\ell(B(x,r))\leq\Delta^{-\sigma}r^a
 \qquad(\Delta\leq r\leq1).
\]
If $0<\Delta<\Delta_0$, then
\begin{equation}
 \left|\bigcup_\ell P_\ell\right|_\Delta
 \geq \Delta^{-F(a,b)+\omega},
 \qquad F(a,b)=a+\min\{1,b,(a+b)/2\}.
 \tag{RW}\label{eq:RW}
\end{equation}
Here $|\cdot|_\Delta$ denotes covering number; the constants may depend on the
fixed charts and on $C$.
\end{theorem}

\begin{proof}
The separated-set version is the point--line dual of the discretized
Furstenberg theorem of Ren and Wang \cite[Theorem 4.1]{RenWang2025}.
We explain the passage from probabilities, which is the only additional
point needed here. Choose $a'<a$ and $b'<b$ sufficiently close to $a,b$
for the prescribed final loss. Take the slack $\sigma$ in the hypothesis
smaller than both $\min(a-a',b-b')/2$ and one quarter of the slack
allowed by the cited separated-set theorem at $a',b'$. Choose the
additional sampling loss below another quarter of that slack. Replace
line parameters by
representatives of their $\Delta$-grid cells. Points incident to a line
remain within $O(\Delta)$ of its representative. On the finite grid sample
$\lceil\Delta^{-b'}\rceil$ line parameters independently from $\nu$.
For a dyadic ball of radius $r\geq\Delta$, the expected number of samples is
at most
\[
 2\Delta^{-\sigma-b'}r^b
 \leq 2\Delta^{-\sigma}(r/\Delta)^{b'}.
\]
For any prescribed small additional power loss, binomial tails bound all
these counts simultaneously by that power times $(r/\Delta)^{b'}$.
Indeed there are only polynomially many grid balls to consider, whereas
the probability of exceeding a sufficiently large small-power multiple
of the displayed expectation and $1$ decays faster than every fixed
power of $\Delta$. This also bounds multiplicities in a single cell.
The probability of one line grid cell is $O(\Delta^{b-\sigma})$, so the
expected fraction of samples repeating an earlier cell is
$O(\Delta^{b-b'-\sigma})$. Discard these repetitions at a small-power
cardinality cost. Perform the same grid sampling on each retained line
with exponent $a'<a$; its point-cell probabilities are
$O(\Delta^{a-\sigma})$. The expected fraction of point repetitions is
$O(\Delta^{a-a'-\sigma})$. The number of lines is polynomial, so the
binomial count bounds hold simultaneously; by averaging, discard any
small fraction of lines with excessive point repetitions. Equalize
cardinalities at a further small-power cost. The resulting sets satisfy
the nonconcentration and cardinality hypotheses of the separated-set
theorem. Its conclusion,
continuity of $F$, and choices of $a',b'$ sufficiently close to $a,b$ give
\eqref{eq:RW}. Small changes of width and projection of the point supports
onto representative lines change covering numbers only by bounded
factors. All small losses are chosen before $\Delta$ tends to zero.
\end{proof}

For a finite or countable discrete name $U$, side information $D$, and
$0<\rho<1$, put
\[
 i_\rho(U\mid D)
 =\frac{-\log\mathbb P(U\mid D)}{\log(1/\rho)},
\]
evaluated at the sampled value of $U$. Conditional probabilities are
measurable versions on the probability spaces of the cylinder construction.

\begin{lemma}[Likelihood and list comparisons]
\label{low:likelihood}
For arbitrary additional side information $D'$ and $\epsilon>0$,
\begin{equation}
 i_\rho(U\mid D)\geq i_\rho(U\mid D,D')-\epsilon
 \tag{LP}\label{eq:LP}
\end{equation}
outside a set of probability at most $\rho^\epsilon$.
If $I(U;D'\mid D)=o(\log(1/\rho))$, the difference of these scores tends
to zero in probability. More precisely, if the sampled $U$ belongs on an
event $E$ to a list $\mathcal L(V,D)$ of at most $\rho^{-b}$ names, then
\[
 \mathbb P\bigl(E\cap\{i_\rho(U\mid D)>
              i_\rho(V\mid D)+b+\epsilon\}\bigr)
 \leq\rho^\epsilon.
\]
If $V$ instead belongs to a list indexed by $(U,D)$, reverse $U,V$ in
this comparison. Thus a change of discrete names by such lists costs at
most $b+\epsilon$ in the indicated direction. The exceptional probability
is measured on $E$ under the original law; no conditioning on $E$ is
implicit. Finally, restricting the law to
an event of probability bounded below preserves any fixed finite collection
of normalized scores up to an error tending to zero in probability under
the restricted law.
\end{lemma}

\begin{proof}
Write $p(u\mid D)$ and $q(u\mid D,D')$ for the two conditional masses.
Conditioning on $(D,D')$ and summing over $u$ gives
\[
 \mathbb P\{q(U\mid D,D')<\rho^\epsilon p(U\mid D)\}
 \leq\rho^\epsilon.
\]
This is \eqref{eq:LP}. If
$Z=\log(q(U\mid D,D')/p(U\mid D))$, the same argument gives
$\mathbb P(Z<-t)\leq e^{-t}$ and hence $\mathbb E Z_-\leq1$.
Since $\mathbb E Z=I(U;D'\mid D)$, Markov's inequality applied to $Z_+$
proves the mutual-information assertion.

For lists, conditional on the known data, the total probability of names
in a list of size $L$ that each have conditional probability less than
$\rho^\epsilon/L$ is at most $\rho^\epsilon$. Combine this observation
with the chain rule for conditional masses and the fact that
$\mathbb P(U\mid D)\geq\mathbb P(U,V\mid D)$ at the sampled names.
The argument still applies if only values attained on the specified
event are listed.

Let $E$ be the restricting event, with $\mathbb P(E)\geq c>0$.
The new conditional mass is
\[
 \mathbb P(U\mid D,E)
 =\mathbb P(U\mid D)
   \frac{\mathbb P(E\mid U,D)}{\mathbb P(E\mid D)}.
\]
Both factors in the ratio are at most $1$. Under the restricted law, the
probability that either is smaller than $\rho^\epsilon$ is at most
$2\rho^\epsilon/c$, by integrating the corresponding conditional
probability of $E$. Thus the normalized logarithm of the ratio tends to
zero in probability. This also proves the assertion with nonatomic side
information.
\end{proof}

\begin{lemma}[Conditional planar incidence bootstrap]
\label{low:conditional-incidence}
Fix $a,g\in(0,1]$ and $\xi>0$. There are a finite collection of pairs
$0<z'\leq z\leq1$ and a tolerance $\chi>0$ with the following property.
Let $(Q,\mathcal L,D)$ vary with $\rho\downarrow0$, where $Q$ and
$\mathcal L$ lie in fixed bounded point and line charts and are incident
up to $O(\rho)$. Subscripts denote grid names of width $\rho^z$.
If, simultaneously on the chosen pairs, with probability tending to one,
\begin{equation}
 i_\rho(Q_{z'}\mid\mathcal L_z,D)\geq az'-\chi,
 \qquad
 i_\rho(\mathcal L_{z'}\mid Q_z,D)\geq gz'-\chi,
 \tag{cond}\label{eq:cond}
\end{equation}
then, with probability tending to one,
\begin{equation}
 i_\rho(Q_1\mid D)\geq a+\min\{1,a+g\}-\xi.
 \tag{inc}\label{eq:inc}
\end{equation}
The same statement holds at any fixed output depth, with the target
multiplied by that depth. The side information and conditional law may
vary along the sequence.
\end{lemma}

\begin{proof}
We first justify one bootstrap step. Suppose the unconditional line
scores, conditional only on $D$, are at least $yz'-\chi$ at a sufficiently
fine finite grid of depths $z'$ between $0$ and a fixed $z>0$.
We claim that the point score at depth $z$ is at least
$(F(a,y)-\omega)z$ in probability, when the input tolerance and mesh are
sufficiently small in terms of $\omega$ and $z$.

Otherwise pass to a subsequence on which violation has probability at
least a fixed $c>0$. A violating point belongs to a cell of conditional
mass greater than $\rho^{(F(a,y)-\omega)z}$; for each $D$ there are fewer
than $\rho^{-(F(a,y)-\omega)z}$ such cells. Intersect this event with all
the good input-score events. Choose a value of $D$ for which this
intersection has conditional probability bounded below, say by $c/2$.
Retain fine line cells on which the conditional probability of the
intersection is at least $c/4$. Their total retained witness mass is
bounded below in terms of $c$. Restrict the original line probability
to these witnesses and normalize. Any occupied coarser line cell
contains a good witness, so its original mass satisfies the required
bound $\rho^{yz'-\chi}$; restriction and normalization cost only a
constant depending on $c$.

In each retained fine line cell, restrict its conditional point
probability to witnesses and normalize. Each occupied coarser point
cell contains a witness to the first inequality of \eqref{eq:cond},
so its mass is bounded by a constant times $\rho^{az'-\chi}$.
Cover Euclidean balls by boundedly many grid cells. Between successive
queried scales use the next larger queried cells; a depth mesh of size
$\vartheta$ costs at most $\rho^{-2\vartheta}$ in the ball estimates.
At the coarsest end the trivial probability bound gives the same
small-power loss. Thus, in units $\Delta=\rho^z$, the retained line
and point measures meet Theorem~\ref{low:ren-wang} with exponents $y,a$
and arbitrarily small nonconcentration slack. Replace each fine line
cell by a representative line; its witnesses move by $O(\Delta)$.
Equation~\eqref{eq:RW}, with a loss less than $\omega z$, contradicts
the upper bound on the number of popular point cells. This proves the
step. Point--line duality proves the symmetric step.

By \eqref{eq:LP}, the starting marginal bounds are $x=a$ and $y=g$,
up to arbitrarily small losses. The two steps iterate
\[
 x\longleftarrow F(a,y),\qquad y\longleftarrow F(g,x).
\]
Use the preceding values on both right-hand sides, so the sequences
are increasing and bounded. At their fixed point, $y>a$ and $x>g$.
For example, $y\leq a$ would give $x=a+y\geq2y$ and then
$F(g,x)>y$, a contradiction. Consequently the fixed-point equations
reduce to
\[
 x=a+\min\{1,(a+y)/2\},\qquad
 y=g+\min\{1,(g+x)/2\}.
\]
This map is a contraction in the maximum norm. Its fixed point is
$(2a+g,a+2g)$ when $a+g\leq1$, and $(1+a,1+g)$ when $a+g\geq1$.
Thus finitely many iterations suffice to reach the strict target in
\eqref{eq:inc}. Choose the finite grids and tolerances backwards from
the final iteration. The argument for output depth $z$ uses
$\rho^z$ as basic scale and is identical after this change of units.
\end{proof}

We now return to a classical cylinder extremal configuration. Assume that
both minimal shear costs vanish, normalize the terminal widths to
$r_1=f_1=1$, and retain the notation of the cylinder construction:
\[
 \Gamma=d+(1-p)/2,\quad p\geq1,\quad 0<d\leq1,
 \qquad E_a=e s^{-da},\quad \kappa_a=\kappa s^{pa}.
\]
The assigned terminal law is $\mathbb P^1$. All following statements use
the ordered limits and finite meshes of that construction. In particular,
ordinary regularizations are inserted in forward order, their total
energies and regularity constants obey the displayed benchmarks up to
$s^{-o(1)}$, the documented terminal count is at most
$s^{-o(1)}E_1/\kappa_1$, and terminal assigned mass is at least
$s^{o(1)}E_1$.

Let $x_c,n_c$ be the unit-bin indices of the exact velocities $X,N$ and set
\[
 A=(X-x_c,N-n_c),\qquad V(A)=(A_1,A_1^2,A_2).
\]
Use boosted spatial coordinates
\[
 (B_t-tx_c,\;Y_t-2x_cB_t+tx_c^2,\;D_t-tn_c).
\]
Their velocity is $V(A)$, and $A$ lies in a fixed bounded square.
For $0\leq a\leq1$, define the nested name
\[
 C(a)=\bigl(x_c,n_c,T_1^{(a)},\operatorname{position}(T_1)^{(a)}\bigr),
\]
where the superscript denotes absolute dyadic bins of width comparable
to $s^a$. Conditional on $C(a)$, refining to $C(b)$ uses at most
$O(s^{-4(b-a)})$ names for $b\geq a$.

\begin{lemma}[Geometric and probabilistic bounds for the boosted names]
\label{low:classical-names}
Suppose $\Gamma>\eta/2$. The preceding configuration has the following
properties, simultaneously on every fixed finite list of depth queries.
\begin{enumerate}
\item A terminal label determines $C(1)$ to $s^{-o(1)}$ possibilities.
After deterministic flattening of the relative scores of $C(a)$ given
$C(0)$, their limiting profile $D(a)$ is nondecreasing and Lipschitz,
with $D(0)=0$ and $D(1)\leq p+d$.
\item For $0<v<1$ and $0\leq a'\leq1$, the typical score of
$A^{(v)}$ given either its terminal label or $C(a')$ is at least
$(3-p)v$ in $\log(1/s)$ units. Every fixed strict deficit has power-small
probability.
\item For fixed $0<a<a+zh\leq1$, the time name at width $s^{a+zh}$,
conditional on the exact particle and $C(a)$, has typical score at least
$dzh$ in $\log(1/s)$ units, with the same meaning of strict deficit.
\item Fix $K=3$, $a>0$, and $h>0$ with $a+Kh<1$. For every
$\epsilon>0$, uniformly over measurable choices of an affine plane
$\Pi_{C(a+Kh)}\subset\mathbb R^3$,
\[
 \mathbb P^1\{\operatorname{dist}(V(A),\Pi_{C(a+Kh)})
                  \leq s^{\epsilon h}\}
 \leq s^{c+o(1)}
\]
for some $c>0$ depending on the fixed parameters and strict gap.
\end{enumerate}
\end{lemma}

\begin{proof}
Zero shear cost and Taylor expansion on a terminal test imply unsheared
velocity widths $s^{-o(1)}$ and spatial widths $s^{1-o(1)}$ after a
bounded velocity boost. The possible unit boost bins therefore cost only
$s^{-o(1)}$ per label. This proves the first assertion about $C(1)$.
Conversely, for each fixed value $c_0$ of $C(0)$, take the union of the
root predecessors of \emph{all} terminal samples with $C(0)=c_0$.
This whole union lies in only boundedly many root time/test pairs of
bounded weight: $A$ and terminal time are bounded, so exact backflow
places the initial boosted positions in a bounded box. Consequently a
list of at most $s^\tau e/\kappa$ values of $C(0)$ has predecessor mass
at most $Cs^\tau e$, for every fixed $\tau>0$. Rarity transfer
\eqref{eq:Sp} makes the corresponding terminal event power-small.
Apply this to the list of atoms with probability greater than
$s^{-\tau}\kappa/e$. This is the popular-list argument of
Lemma~\ref{cyl:phase-name}, and gives the typical lower score
$\log_{1/s}(e/\kappa)-o(1)$ for $C(0)$ without counting terminal times
separately. The documented count
and the label-to-name list give the typical upper score
$\log_{1/s}(E_1/\kappa_1)+o(1)$ for $C(1)$. Their difference is $p+d$.
Apply forward deterministic flattening at final readout on successively
finer finite lists. The relative branching bound and
Lemma~\ref{low:likelihood} imply that all limiting increments lie between
$0$ and $4(b-a)$. Dense-list extension gives the asserted profile and
endpoint bound.

For the angular assertion, ordinary-regularize at $j=1-v$. Fix a
terminal label and an angular $s^v$-cell. In duration units of the
$j$-interval, terminal spatial uncertainty is $s^{v-o(1)}$ in each of
the three boosted spatial coordinates. Backflow through a duration at
most $1$ adds uncertainty $O(s^v)$ horizontally and in the normal
position. After tilting by the angular-cell center, the velocity
curvature contributes $O(s^{2v})$ in the last base coordinate, while
terminal position uncertainty there is still $s^{v-o(1)}$.
Consequently $s^{-v-o(1)}$ full tests with $r=f=s^{v-o(1)}$ cover all
predecessors. Each has weight $s^{4v-o(1)}$, so their total mass is at
most
\[
 s^{-o(1)}\kappa_j s^{3v}
 =s^{-o(1)}\kappa_1s^{(3-p)v}.
\]
Classical masks are pointwise contractions and do not enlarge this
mass. Labels of read mass below $s^\tau\kappa_1$ have total probability
at most $s^{\tau-o(1)}$ by the count. On the remaining labels the
conditional angular atom bound follows, with arbitrarily small $\tau$.
A label predicts $C(a')$ with a subpower list. Testing popular angular
atoms given $C(a')$, and pulling these lists back to labels, gives the
same bound by Lemma~\ref{low:likelihood}.

For the time assertion apply Lemma~\ref{cyl:time-prediction}, with the
exact particle as common side information and $C(a)$ as target. At each
of the cuts $a$ and $a+zh$, use the current time and boosted-position bins
at the common width $s^a$ as relay names. Bounded boosted velocity moves
position by only $O(s^a)$ throughout the depth-$a$ interval. Thus the
earlier row predicts the later relay with bounded ambiguity, and that
relay predicts $C(a)$ with bounded ambiguity. The lemma and \eqref{eq:Sp}
exclude lists of fewer than $s^{-dzh+\tau}$ fine time names for any
fixed $\tau>0$. Testing the atoms that would violate the stated score
gives the assertion. Approximating depths on the finite cut mesh costs
only the ordered vanishing exponent errors.

It remains to prove plane avoidance. Write the plane condition as
\[
 |\alpha A_1+\beta A_1^2+\gamma A_2+c_0|
 \lesssim s^{\epsilon h},\qquad
 \alpha^2+\beta^2+\gamma^2=1.
\]
Choose a fixed $h'>0$ sufficiently small compared with $\epsilon h$,
with $i=1-h'>a+Kh$, and put $\Delta=s^{h'}$.
Fix a small $\lambda>0$. When $|\gamma|\leq\Delta^\lambda$, the
boundedness of $A_2$ reduces the condition to a quadratic sublevel set
in $A_1$ of thickness $O(\Delta^\lambda)$. At least one of
$|\alpha|,|\beta|$ is bounded below. Such a sublevel set is covered by
boundedly many intervals of length $O(\Delta^{\lambda/2})$: when the
quadratic coefficient is bounded below this follows by factoring at
the vertex, and otherwise the derivative is bounded below on the
bounded interval, after a bounded partition. A terminal label knows
only subpower many choices of the conditioning chart. Apply the
zero-cost case of \eqref{eq:mem} with $L=1$ on a final depth gap
$\lambda h'/3$. Its cell width is larger than these intervals, so a
bounded list suffices for each chart. The exponent
$1-d+2\Gamma$ is positive. This part has power-small probability.

Suppose now $|\gamma|>\Delta^\lambda$. Ordinary-regularize at $i$.
A row at this cut knows only subpower many compatible boost and
$C(a+Kh)$ charts, since its exact particle and current time start
predict the coarser name. Make this subpower number of copies and
restrict each to the selected plane slab. In its segment duration
units subtract the velocity shear
\[
 N_{\mathrm{sh}}(A_1)
 =-\gamma^{-1}(c_0+\alpha A_1+\beta A_1^2).
\]
The corresponding spatial shear is one of the allowed cylinder
shears. Its coefficients are $O(\Delta^{-\lambda})$ on any nonempty
bounded chart. Choose $h'$ small enough that
$s^{\epsilon h}\Delta^{-\lambda}\leq\Delta$.
The residual velocity therefore has magnitude $O(\Delta)$.
Partition the residual initial normal position into intervals of
length $\Delta$. On each such slab the base marginal satisfies
\[
 \mu\{|X-x|\leq R,\ |B-b|\leq R,
             |Y-y-2x(B-b)|\leq R^2\}
 \lesssim s^{-o(1)}\kappa_i R^{3+\eta}\Delta^{1-\eta},
 \qquad 0<R\leq1.
\]
Indeed the specified set is contained in a full input test with
normal width comparable to $\Delta$. Since $R^{3+\eta}\leq R^3$,
\eqref{eq:kin} applies with constant
$s^{-o(1)}\kappa_i\Delta^{1-\eta}$.

For each chart and initial normal slab, let $e_*$ be its input mass
and $m_{t,b}$ its base position-cell masses on the $\Delta$ time grid.
Equation~\eqref{eq:kin} and the trivial time-copy bound give
\[
 \sum_{t,b}m_{t,b}^2
 \lesssim s^{-o(1)}\kappa_i\Delta^{2-\eta}e_*,
 \qquad \sum_{t,b}m_{t,b}\lesssim\Delta^{-1}e_*.
\]
Cauchy--Schwarz therefore bounds their $3/2$-sum by
$s^{-o(1)}\sqrt{\kappa_i}\Delta^{(1-\eta)/2}e_*$.
Subsequent classical contractions only decrease these masses.
A terminal label meets at most
$s^{-o(1)}\Delta^{-O(\lambda)}$ base position cells and initial
normal slabs: its spatial diameter after the shear is
$s^{-o(1)}\Delta^{1-O(\lambda)}$, and residual velocity is
$O(\Delta)$ over the unit duration. Each label also sees only
subpower many charts. Summing the preceding estimate and using
the finite-list convexity inequality for each label yields
\[
 \sum_\nu m_\nu^{3/2}
 \lesssim s^{-o(1)}\Delta^{(1-\eta)/2-O(\lambda)}
                  E_i\sqrt{\kappa_i},
\]
where $m_\nu$ are terminal masses restricted to the plane event.
Their total is at least
$s^{o(1)}\mathbb P^1(\text{event})E_1$. The terminal count and
H\"older's inequality yield the reverse estimate
\[
 \sum_\nu m_\nu^{3/2}
 \geq s^{o(1)}\mathbb P^1(\text{event})^{3/2}
          E_1\sqrt{\kappa_1}
 =s^{o(1)}\mathbb P^1(\text{event})^{3/2}
   \Delta^{1/2-\Gamma}E_i\sqrt{\kappa_i}.
\]
Thus the event probability to the power $3/2$ is at most
$s^{-o(1)}\Delta^{\Gamma-\eta/2-O(\lambda)}$.
Choose $\lambda$ sufficiently small relative to $\Gamma-\eta/2$.
This proves the fourth assertion, uniformly in the plane choice.
Maximizing conditionally over a countable sufficiently fine parameter
grid, and then enlarging the slab by a bounded factor, also shows that
the conditional supremum of slab probabilities tends to zero for
typical $C(a+Kh)$ names. Measurable approximate maximizers suffice.
\end{proof}

\begin{proposition}[Exclusion of classical zero costs]
\label{low:classical-zero-cost}
A classical cylinder extremal configuration with
$\Gamma>\eta/2$ cannot have both minimal shear costs equal to zero.
\end{proposition}

\begin{proof}
Suppose such a configuration exists and use
Lemma~\ref{low:classical-names}. Put $m=3-p$.
The relation between $p,d,\Gamma$ gives
$m=2-2d+2\Gamma>0$, while $p\geq1$ gives $m\leq2$.
The profile supplied by that lemma satisfies $D(1)\leq p+d$.
We will prove $D'(a)\geq1+3d$ at almost every interior depth, contradicting
$p+d=1+3d-2\Gamma$. At a short gap, conditional sampling and planar
incidences will force information in a two-dimensional projection of
the spacetime position. Recovering the two omitted coordinates from
time uncertainty gives the required derivative bound.

We first identify depth parameters at which conditioning changes angular
information negligibly on a short depth interval. For each dyadic $h>0$,
take subsequential limits, simultaneously on dense depth lists, of
\[
 f_h(a)=\lim\frac{H(A^{(Kh)}\mid C(a))}{\log(1/s)}.
\]
These functions have values in $[0,2Kh]$, are nonincreasing, and are
Lipschitz with constant $4$, by relative branching of $C$.
The limiting normalized information equals
$(f_h(a)-f_h(a+Kh))/h$. Integrating it over $0<a<1-Kh$ gives at most
$2K^2h$. Summing over dyadic $h$ and applying Tonelli's theorem shows
that, for almost every interior $a$,
\[
 \frac{I(A^{(Kh)};C(a+Kh)\mid C(a))}{h\log(1/s)}
 \longrightarrow0
 \quad\text{as dyadic }h\downarrow0
\]
in the ordered limiting sense. Fix such an $a$ at which $D$ is also
differentiable.

Write $\rho=s^h$, $C=C(a)$, and $C'=C(a+Kh)$.
Subtract the lower endpoints of the time and position bins of $C$
and divide by their common width. The finer name $C'$ determines a
representative $O\in\mathbb R^4$ with error $O(\rho^K)$ for
\[
 (0,W)+\theta e(A),\qquad e(A)=(1,A_1,A_1^2,A_2).
\]
Here $W$ is determined by the exact particle and $C$, by evaluating
position at the start of the time bin; all displayed variables are
bounded. At each fixed finite collection of queried $z>0$,
Lemma~\ref{low:classical-names} gives the lower scores
\[
 i_\rho(A_z\mid C')\geq mz-o(1),\qquad
 i_\rho(\theta_z\mid C,\text{particle})\geq dz-o(1)
\]
in probability.

Draw $M$ rows independently conditional on $C'$, with $M$ any fixed
integer, and put $L_i=A_i^{(Kh)}$, using bin representatives when
evaluating $e_i=e(L_i)$. The rows share $O$ and $C$.
The time bound survives conditioning on all pins. Indeed $L_i$ is
known from particle $i$, and conditional independence gives
\[
 I(\mathbf L_{\ne i};\text{particle}_i,\theta_i\mid C)
 \leq I(\mathbf L_{\ne i};C'\mid C)
 \leq\sum_{j\ne i}I(L_j;C'\mid C)
 =o(\log(1/\rho)).
\]
For completeness, the second inequality follows by expanding conditional
entropy: the pins have additive entropy given $C'$, whereas their
entropy given $C$ is at most the sum of their individual entropies.
The first inequality is conditional data processing. The information
bound also controls the information about $\theta_i$ conditional on
its particle, by the mutual-information chain rule.
Apply Lemma~\ref{low:likelihood}. Independently, the angular bounds
conditional on $C'$ survive conditioning on the other pins by the
product law.

For every fixed small $\epsilon>0$, with probability tending to one,
all four-element subsets satisfy
\[
 |\det(e_i,e_j,e_k,e_\ell)|\geq\rho^{O(\epsilon)},
 \qquad |(L_i)_1-(L_j)_1|\geq\rho^{O(\epsilon)}\quad(i\ne j).
\]
To see this, sample the vectors successively. Uniform conditional plane
avoidance from Lemma~\ref{low:classical-names} excludes a slab about
any plane containing the affine span of previously sampled $V$'s.
The resulting lower bounds for successive distances multiply to a
determinant bound. For fewer than three predecessors extend their span
to a plane. Vertical planes give the second inequality. A union bound
suffices because $M$ is fixed. Rounding errors $O(\rho^K)$ are smaller
than all these tolerances. The implied power constants are absolute.

For $i<j$, let $\pi_{ij}$ be orthogonal projection to the quotient
plane with kernel $\operatorname{span}(e_i,e_j)$, using measurable
orthonormal coordinates. Write $g_{ij}(A_k)$ for the projective
direction of $\pi_{ij}e(A_k)$. Choose the finite query depths that
will be needed in Lemma~\ref{low:conditional-incidence}. Color each
increasing triple $i<j<k$ by the vector of $\epsilon$-quantized scores
\[
 i_\rho\bigl((g_{ij}(A_k))_z\mid C',\mathbf L_{\ne k}\bigr)
\]
at these depths. There are finitely many colors: scores outside
$[0,2]$ can be assigned one overflow color, which has probability
tending to zero because a direction has $O(\rho^{-z})$ cells and
$z\leq1$. Take $M$ sufficiently large that the finite Ramsey theorem
\cite[Theorem~B, p.~267]{Ramsey1930}
provides four indices all of whose triples have the same color.
Only this monochromatic four-set conclusion for a finite coloring of
triples is used here.

We verify the recovery estimate used with this monochromatic set.
For its largest index $k$, choose the smallest index $i$ and the other
two indices $j,\ell$. A direction for $\pi_{ij}e(A_k)$ specifies the
three-dimensional linear space
$\operatorname{span}(e_i,e_j,e(A_k))$; the other direction specifies
$\operatorname{span}(e_i,e_\ell,e(A_k))$.
The determinant bound implies that these spaces meet transversely in
$\operatorname{span}(e_i,e(A_k))$, with inverse norms at most
$\rho^{-O(\epsilon)}$. Thus two direction cells of width $\rho^z$
constrain any bounded candidate with first coordinate $1$ to within
$O(\rho^{z-O(\epsilon)})$ of the affine line through $e_i,e(A_k)$.
On this line, parameterized by $(1-t)e_i+te(A_k)$, the equation
``third coordinate equals the square of the second coordinate''
has residual
\[
 t(1-t)\bigl((A_k)_1-(L_i)_1\bigr)^2.
\]
The roots $0,1$ are simple with derivative bounded below by a power
$\rho^{O(\epsilon)}$. Consequently the candidate angular parameter
lies in a bounded number of balls of radius
$\rho^{z-O(\epsilon)}$. This is a list of at most
$\rho^{-O(\epsilon)}$ cells at the queried width.

It follows from the angular lower bound and list comparison that the
joint score of these two directions is at least
$mz-O(\epsilon)-o(1)$, on transverse samples outside a vanishing
exceptional set. The sum of their individual scores is at least the
joint score minus $o(1)$ by \eqref{eq:LP} and the chain rule.
Their colors agree, so each individual score is at least
$(m/2)z-O(\epsilon)-o(1)$, simultaneously at all query depths.
Since there are only finitely many index choices, one fixed triple
$i<j<k$ has these properties on an event of probability bounded below
independently of $\rho$. Restrict to that event. By
Lemma~\ref{low:likelihood} all previously obtained finite score bounds
remain valid in probability, with vanishing additional losses.
One can further restrict to a fixed bounded slope chart after one of
finitely many fixed rotations, retaining positive probability.

Set $D_*=(C,\mathbf L_{\ne k})$ and define
\[
 Q=\pi_{ij}O,\qquad
 \mathcal L=\pi_{ij}(0,W_k)+\mathbb R\pi_{ij}e(A_k).
\]
These form a point--line incidence to accuracy $O(\rho^K)$ in bounded
charts. We check \eqref{eq:cond} with $a=d$ and $g=m/2$.
After adjoining particle $k$, the line is known. Its projected speed
is at least $\rho^{O(\epsilon)}$, by transversality, so a point cell
of width $\rho^{z'}$ permits at most $\rho^{-O(\epsilon)}$ time
cells of that width. The conditional time bound and
\eqref{eq:LP} prove the first inequality of \eqref{eq:cond}.
For the second, adjoin $C'$. The point name is then known, and a line
cell specifies its direction to boundedly many direction cells.
The score bound obtained from the monochromatic triple and
\eqref{eq:LP} prove the required line score. All tolerances can be
made smaller than the tolerance in
Lemma~\ref{low:conditional-incidence}. Hence
\[
 i_\rho(Q_1\mid D_*)
 \geq d+\min\{1,d+m/2\}-o(1)
\]
with an arbitrarily small prescribed strict loss.

We must also recover the two coordinates removed by $\pi_{ij}$.
Complete the quotient coordinates to an orthonormal basis using first
the unit vector along the part of $e_j$ perpendicular to $e_i$, and
then the unit vector along $e_i$. Given particle $j$ and $D_*$, the
quotient point $Q_1$ has only boundedly many possibilities: its
velocity component is annihilated up to $O(\rho^K)$, because $L_j$
is known. The next coordinate has speed at least
$\rho^{O(\epsilon)}$ and therefore reveals $\theta_j$ to a list of
$\rho^{-O(\epsilon)}$ cells. Its score conditional on $Q_1,D_*$ is
at least $d-O(\epsilon)$ by the time bound, the bounded-list version
of the chain rule, and \eqref{eq:LP}. For the last coordinate adjoin
particle $i$. Both preceding coordinates are predictable with bounded
ambiguity, whereas the last coordinate reveals $\theta_i$ with the
same accuracy. This contributes another $d-O(\epsilon)$.
The chain rule, bounded changes between orthonormal coordinate grids,
and \eqref{eq:LP} to remove the pins give
\[
 i_\rho(O_1\mid C)
 \geq 3d+\min\{1,d+m/2\}-\text{arbitrarily small strict loss}.
\]

The order of choices is as follows. First choose the target strict
loss and the finite incidence iterations and query grids. Next choose
the color accuracy, transversality tolerance, and input score
tolerances sufficiently small. Then fix a finite Ramsey number $M$.
Only after these choices let dyadic $h\downarrow0$ at the selected
$a$, taking all preceding scale and mesh limits first for each $h$.
The stationary information loss then tends to zero, and every event
restriction after this information comparison has probability bounded
below by a constant depending on $M$, not on the inner scale.

Given $C$, the name $O_1$ and the refinement $C(a+h)$ determine each
other with bounded ambiguity. Thus its limiting score in
$\log(1/\rho)$ units is
$(D(a+h)-D(a))/h$, which tends to $D'(a)$. Since
$d+m/2=1+\Gamma>1$, the preceding inequality gives
$D'(a)\geq1+3d$ for almost every $a$. Integrating and using
Lemma~\ref{low:classical-names} gives
\[
 1+3d\leq D(1)\leq p+d=1+3d-2\Gamma,
\]
a contradiction. This proves the proposition.
\end{proof}
